% Introdução do artigo (template SBC) — Lab03S01, Issue #15.
% As hipóteses abaixo foram escritas ANTES de qualquer coleta ou análise de dados do Lab03
% (o commit que adiciona este arquivo precede a execução do pipeline para os 100 repositórios).
\section{Introdução}

As métricas DORA (\textit{DevOps Research and Assessment}) tornaram-se o padrão de mercado para medir
o desempenho de entrega de software~\cite{forsgren2018accelerate}. Quatro medidas descrevem duas
dimensões: velocidade (\textit{deployment frequency} e \textit{lead time for changes}) e estabilidade
(\textit{change failure rate} e tempo de recuperação de um deploy que falhou). O relatório
\textit{Accelerate State of DevOps} sustenta que as equipes de melhor desempenho vão bem nas duas
dimensões, ou seja, que velocidade e estabilidade não são um \textit{trade-off}.

O GitHub, porém, não registra ``deploys em produção'' nem ``falhas em produção''. Ao minerar
repositórios open-source, só é possível \emph{inferir} as métricas por \textbf{proxies}: uma
\textit{release} publicada faz o papel de deploy, e uma execução de CI que falhou faz o papel de falha.
Proxies podem enganar. Uma biblioteca que publica uma release por mês não ``faz deploy'' em produção;
quem faz isso são os usuários dela. Por isso o valor de uma métrica minerada precisa vir acompanhado
de uma medida de confiança.

Este trabalho tem como \textbf{objetivo} calcular automaticamente as métricas DORA para repositórios
populares que usam GitHub Actions e medir o quanto as conclusões dependem das definições operacionais
escolhidas. A seguir, uma hipótese informal por questão de pesquisa, escrita antes de ver os dados.

\begin{description}
  \item[H1 (RQ01, frequência de deploy).] Esperamos uma distribuição muito assimétrica, com mediana
  entre uma release a cada duas semanas e uma por mês (categoria \textit{Medium}) e uma cauda de
  projetos com mais de uma release por semana. Projetos populares costumam seguir um ritmo regular de
  versões, mas poucos publicam diariamente.

  \item[H2 (RQ02, \textit{lead time}).] A mediana por commit (variante b) será bem menor que a
  mediana por release (variante a). Um único commit antigo incluído na release faz a variante (a)
  crescer, enquanto a (b) é pouco afetada. Esperamos lead times medianos de semanas a poucos meses,
  pois releases agrupam muitos commits acumulados.

  \item[H3 (RQ03, taxa de falha).] A taxa de falha de CI (variante a) será maior que a taxa de
  releases corretivas (variante b). Falhas de pipeline são comuns (testes instáveis, lint) e raramente
  chegam a uma release corretiva. Esperamos um CFR (a) mediano entre 10\% e 30\%.

  \item[H4 (RQ04, tempo de recuperação).] A maior parte dos episódios de falha se resolve em horas
  (correção rápida no mesmo dia), mas a cauda é longa: dias ou semanas em projetos com pouca
  manutenção. A mediana por repositório ficará entre 1 hora e 1 dia, e uma parcela relevante dos
  episódios será censurada.

  \item[H5 (RQ05, frequência $\times$ falha).] Esperamos uma correlação de Spearman fraca ou
  levemente negativa entre frequência de deploy e CFR. Isso seguiria a afirmação do DORA de que não há
  \textit{trade-off}, mas com ruído grande por causa dos proxies e da diferença entre tipos de projeto.

  \item[H6 (RQ06, características dos repositórios).] A linguagem principal e o tipo de projeto
  (biblioteca, aplicação, ferramenta CLI) terão associação mais forte com as métricas do que a
  popularidade. Aplicações e serviços devem publicar com mais frequência que bibliotecas. Idade e
  número de contribuidores terão efeitos pequenos, e a maior parte das diferenças perderá significância
  após a correção para comparações múltiplas.

  \item[H7 (RQ07, sensibilidade).] A classificação DORA mudará de categoria para uma fração
  considerável dos repositórios (esperamos mais de 30\%) quando trocamos a definição de deploy
  (release, release com pré-release ou tag) e a variante de \textit{lead time} e de CFR. A
  concordância ponderada (kappa) entre combinações será no máximo moderada, o que indica que parte das
  conclusões das demais RQs depende de escolhas arbitrárias.
\end{description}
